\subsection{两个全嵌入}\label{sec:18f}
此时，引入一个技术条件很方便。若在消失路径特选基的每个同伦类中，都有一个基使其相应消失圈处于一般位置，就称 $\pi:E\to S$ 具有\emph{一般平行移动}。给定任意 Lefschetz 纤维化，都能找到 $(\theta_E,\omega_E,I_E)$ 的一个小扰动，使其在 $\partial E$ 及 $M=E_*$ 附近恒定，且扰动后的结构具有一般平行移动。这很容易从\stref{引理}{15.3}及特选基的同伦类只有可数多个这一事实推出。这种扰动不改变 $M$ 或 $E$ 的 Fukaya 范畴；后一情形指在拟同构意义下不变，并使用 Moser 引理。因此，从现在起假设给定的 Lefschetz 纤维化已具有一般平行移动。

从现在起要求 $\operatorname{char}(\K)\ne2$。这使我们能应用第\sr{sec:6}{6 节}的结果；更重要的是，它使\stref{引理}{18.9}中的 $\Id$ 与 $-\Id$ 不同。\znote{282}{这个假设不是解析上的方便，而是构造方式的副产品：全节把 $\widetilde E$ 当作 $E$ 的二重分歧覆盖来用，而对 $\Z/2$ 取不变部分（如引理 18.9(i) 中对 $HF^*(\widetilde\Delta^\#,\widetilde\Delta^\#)$ 取不变部分）必须能区分 $\pm1$ 两个本征值；特征 $2$ 时二者合并，信息全失。§18j 末尾明说这是为整套论证付出的「唯一代价」（印刷页 288）。因此引用定理 18.24、推论 18.25、18.26 时，不能把结论直接搬到 $\operatorname{char}\K=2$ 的情形；§6 的一般代数结果同样以此为前提。[来源：原书 §18f、§18j（印刷页 282、288）]}

取 $\widetilde{\A}\subset\F(\widetilde E)^{\mathrm{naive}}$ 为全 $\Ai$ 子范畴，其对象为预等变膜 $\widetilde L_\pm^\#$ 及等变膜 $\widetilde L^\#,\widetilde\Delta_\gamma^\#$。虽然下面大体上是在重复前面的讨论，我们仍想明确说明构造 $\widetilde{\A}$ 时使用哪些 Floer 数据与扰动数据。对于任意一对 (U) 型等变膜或预等变膜，应选择使\eref{18.20}成立的 Floer 数据。特别地，这蕴含每个这样的膜都是第\sr{sec:6a}{6a 小节}意义下 $H^0(\widetilde{\A})$ 的自由对象。类似地，对于 (U) 型膜的任意 $(d+1)$ 元组，$d\ge2$，取同类扰动数据，使该同构与 $\Ai$ 结构映射相容。其次，对于任意一对 (B) 型等变膜，Floer 数据应为满足\stref{假设}{14.5}的不变数据。不过，在这些膜来自满足\eref{18.22}的路径 $\gamma_0,\gamma_1$ 的特殊情形中，要采用导出同构\eref{18.23}的更受限制的选择。每当消失路径集 $(\gamma_0,\ldots,\gamma_d)$ 满足\eref{18.24}时，对扰动数据的选择也作同样要求。这与整体选择的一致性条件相容，因为这类 $(d+1)$ 元组对于取有序子集封闭。

\SourcePage{290}
\StatementNumber{12}\begin{definition}\label{st:18.12}
\emph{Lefschetz 纤维化的 Fukaya 范畴}记为 $\F(\pi)$，定义为不变部分 $\widetilde{\A}^{\mathrm{inv}}$。
\end{definition}

\StatementNumber{13}\begin{proposition}\label{st:18.13}
存在在上同调上全忠实的 $\Ai$ 函子
\[
\F(E)\longrightarrow\F(\pi),
\]
将每个位于 $E^{\mathrm{shrink}}$ 中的正合 Lagrangian 膜 $L^\#$ 映到其原像 $\widetilde L^\#$。
\end{proposition}
\begin{proof}
由假设，$E$ 中每个正合 Lagrangian 膜都同痕于一个位于 $E^{\mathrm{shrink}}$ 中的膜。因此，由这类膜构成的全 $\Ai$ 子范畴到 $\F(E)$ 的包含是拟等价，由\stref{定理}{2.9}可将其反演。事实上，考察该定理前的论证，就能找到一个逆拟等价，使它在该子范畴自身上的限制为恒等。另一方面，由\eref{18.20}，同一子范畴可等同于 $\widetilde{\A}^{\mathrm{inv}}$ 中由对象 $\widetilde L^\#$ 构成的全 $\Ai$ 子范畴。
\end{proof}

现令 $\symbfit{\gamma}=(\gamma_1,\ldots,\gamma_m)$ 为消失路径的特选基。对每个 $k=1,\ldots,m$，选择一个分次，使相关的 (B) 型 Lagrangian 子流形 $\widetilde\Delta_k$ 成为等变膜，并给相应消失圈 $V_k$ 赋予诱导的膜结构。分别以 $\widetilde{\symbfit{\Delta}}^{\#}$ 和 $\symbfit V^\#$ 表示由此得到的 $\F(\pi)$ 与 $\F(M)$ 的对象集。它们带有相应的有向 $\Ai$ 子范畴 $\F^\to(\widetilde{\symbfit{\Delta}}^{\#})$ 与 $\F^\to(\symbfit V^\#)$。

\StatementNumber{14}\begin{proposition}\label{st:18.14}
$\widetilde{\symbfit{\Delta}}^{\#}$ 是 $\F(\pi)$ 中的例外集。此外，$\F^\to(\widetilde{\symbfit{\Delta}}^{\#})$ 与 $\F^\to(\symbfit V^\#)$ 拟同构，拟同构由保持对象次序的严格含幺 $\Ai$ 函子给出。因此，存在在上同调上全忠实的 $\Ai$ 函子
\[
\F^\to(\symbfit V^\#)\longrightarrow\F(\pi),
\]
将每个 $V_k^\#$ 映到 $\widetilde\Delta_k^\#$。
\end{proposition}
\begin{proof}
先假设消失圈处于一般位置。由\stref{引理}{18.9}(i)，有
\[
\Hom_{H\F(\pi)}(\widetilde\Delta_k^\#,\widetilde\Delta_k^\#)
=HF^*(\widetilde\Delta_k^\#,\widetilde\Delta_k^\#)^{\Z/2}
\cong H^*(S^{n+1};\K)^{\Z/2}=\K,
\]
故每个 $\widetilde\Delta_k^\#$ 都是例外对象。此外，如\stref{笔记}{18.10}中所述，同一引理的 (ii) 及\eref{18.23}说明，$j>k$ 时，$HF^*(\widetilde\Delta_j^\#,\widetilde\Delta_k^\#)$ 上的 $\Z/2$ 作用为 $-\Id$。因此 $\F(\pi)$ 中的相应态射群无上同调，说明这些对象确实构成例外集。\stref{引理}{5.20}于是给出从 $\F^\to(\widetilde{\symbfit{\Delta}}^{\#})$ 到 $\F(\pi)$ 的在上同调上全忠实函子。另一方面，由\eref{18.23}，该有向子范畴与 $\F^\to(\symbfit V^\#)$ 相同。

\SourcePage{291}
在一般情形中，使用 Lefschetz 纤维化具有一般平行移动这一假设。它说明，对任意 $\symbfit{\gamma}$，都有一个与之同伦的 $\symbfit{\gamma}'$，使上述论证适用。每个消失圈 $V'_k$ 都同痕于 $V_k$，故它们成为 $H^0\F(M)$ 的同构对象。由\stref{引理}{5.21}，所得两个有向 $\Ai$ 子范畴拟同构。同样，每个 $\widetilde\Delta'_k$ 都通过 $\Z/2$-不变 Lagrangian 子流形同痕于 $\widetilde\Delta_k$，因而 $H^0\F(\pi)$ 中的相应对象同构。这就得到相应有向子范畴之间，以及全 $\Ai$ 子范畴之间的拟同构。
\end{proof}
